%! Right Angles Revisited --- Setting Up Orthogonality — Unit 4, Lesson 4.0 — Warm-Up
\documentclass[10pt]{article}
\usepackage{linalg-article}
\usepackage{linalg-boxes}
\newcommand{\TallMath}[1]{$\displaystyle #1\rule[-1.4em]{0pt}{3.2em}$}
\newcommand{\vv}{\mathbf{v}}
\newcommand{\ww}{\mathbf{w}}

\begin{document}

\pageheader{Unit 4, Lesson 4.0}{Warm-Up}
\namedateperiod

\vspace{0.12in}

\begin{notesbox}{Getting Ready: Skills We'll Use Today}
\small These three Unit~1 skills (Lessons 1.0 and 1.2) power all of Unit~4. Answer quickly.

\vspace{4pt}
\textbf{1. Dot product.} With $\vv=\begin{bmatrix} 2 \\ 1 \end{bmatrix}$ and
$\ww=\begin{bmatrix} 3 \\ 4 \end{bmatrix}$, multiply matching components and add:
\[
  \vv\cdot\ww = (2)(3) + (1)(4) = \blank{2.2cm}.
\]

\vspace{4pt}
\textbf{2. Length and unit vector.} For $\vv=\begin{bmatrix} 3 \\ 4 \end{bmatrix}$:
\[
  \|\vv\| = \sqrt{\vv\cdot\vv} = \sqrt{3^2+4^2} = \blank{1.4cm},
  \qquad
  \frac{\vv}{\|\vv\|} = \begin{bmatrix}\rule{0pt}{1.1em}\blank{1.1cm}\\[2pt]\blank{1.1cm}\end{bmatrix}.
\]

\vspace{4pt}
\textbf{3. A right angle?} Two vectors are \emph{perpendicular} exactly when their dot product is
$0$. Are $\begin{bmatrix} 3 \\ 4 \end{bmatrix}$ and $\begin{bmatrix} 4 \\ -3 \end{bmatrix}$
perpendicular? Compute the dot product and decide.
\writeline
\end{notesbox}

\end{document}
